\section{HPS Apparatus, Data-Taking Campaigns, and Analysis Scope}
\label{sec:datasets}

\subsection{Detector configurations and apparatus relevant to the prompt search}

HPS is a forward spectrometer optimized for narrow $e^+e^-$ resonances and displaced
vertices in the high-rate environment immediately downstream of a thin tungsten target.
Its prompt $A^\prime \rightarrow e^+e^-$ sensitivity is set primarily by three
subsystems: the Silicon Vertex Tracker (SVT), which determines the mass and vertex
resolution; the electromagnetic calorimeter (ECal), which provides the trigger and
track--cluster matching; and the beamline/target geometry, which defines the forward
acceptance and the occupancy environment \cite{HPSExperiment2022,HPS2016PromptLong}.

The 2015 and 2016 engineering runs used the baseline six-station SVT. In that
configuration, six double-layer stereo measurement stations were distributed between
$z=\SIrange{10}{90}{cm}$ downstream of the target, with the first active silicon only
about \SI{1.5}{mm} from the beam plane \cite{HPS2016BumpInternal,HPSExperiment2022}.
The ECal geometry used in the engineering runs was already close to the final HPS
layout: two PbWO$_4$ crystal halves with the highest-rate crystals nearest the beam
removed, leaving the beam to pass between the two halves
\cite{HPS2016BumpInternal,HPSExperiment2022}.

Before the first physics run in 2019, HPS upgraded the front tracker. The 2019 and
2021 running periods used a seventh double layer at $z=\SI{5}{cm}$, thin silicon,
slim-edge processing, and short split striplets to preserve low-mass acceptance while
reducing material and occupancy \cite{HPSExperiment2022,HPSSlimEdgeSensors2020}. The
same slim-edge module design was used to replace the previous first layer, and the
front layers were moved closer to the beam. Those hardware changes are the main reason
the 2019/2021 configuration improves the low-mass prompt search: the first precision
measurements occur closer to the beam, the front material budget is reduced, and the
small-opening-angle acceptance is correspondingly better. The ECal crystal layout
itself did not change between the engineering and physics runs, but a positron
hodoscope was installed in front of the ECal before the 2019 run to support the
single-arm positron trigger that recovers events otherwise lost in the ECal hole
\cite{HPSExperiment2022}. Figure~\ref{fig:hps-svt-eras} shows the baseline and
upgraded SVT layouts side by side because that detector evolution is what later drives
the differences in mass resolution and low-mass prompt acceptance used throughout the
analysis note.

\begin{figure}[H]
\centering
\begin{subfigure}[t]{0.90\linewidth}
  \centering
  \graphicorplaceholder{\linewidth}{apparatus_figs/hps_2015_2016_svt_baseline.png}
  \caption{Baseline six-layer SVT used in the 2015 and 2016 engineering runs
  \cite{HPS2016BumpInternal}.}
\end{subfigure}

\vspace{0.9em}

\begin{subfigure}[t]{0.78\linewidth}
  \centering
  \graphicorplaceholder{\linewidth}{apparatus_figs/hps_2019_2021_svt_upgrade.png}
  \caption{Front-layer upgrade used in the 2019 and 2021 physics runs, including the
  added seventh double layer and slim-edge front sensors
  \cite{HPSExperiment2022,HPSSlimEdgeSensors2020}.}
\end{subfigure}
\caption{SVT configurations relevant for the prompt-resonance note. The 2015/2016
baseline tracker and the 2019/2021 upgraded front tracker differ in both layer count
and the construction of the innermost stations, while the ECal remains the same
subsystem apart from the addition of the positron hodoscope before the physics runs.}
\label{fig:hps-svt-eras}
\end{figure}

\subsection{Data-taking campaigns and present analysis scope}

The present note is organized around three HPS running campaigns: the 2015 engineering
run, the 2016 engineering run, and the 2021 physics run. The 2015 and 2016 datasets
are especially natural to compare because they were recorded with the same baseline
detector configuration and were already analyzed in the published prompt-search
workflow. The 2021 dataset belongs to the upgraded detector configuration and is the
first large-statistics physics-run sample to be carried through the GPR workflow.

The 2019 physics run is acknowledged here because it shares the upgraded SVT and
hodoscope configuration with 2021, but its prompt-resonance analysis is not included in
the present scope. This note keeps the statement explicit: 2019 data exist, but they
will be analyzed in a similar capacity in later work rather than merged into the
review-ready workflow described here.

\begin{table}[H]
\centering
\small
{\setlength{\tabcolsep}{6pt}
\begin{tabularx}{\linewidth}{@{} P{0.24\linewidth} P{0.12\linewidth} P{0.16\linewidth} P{0.23\linewidth} Y @{}}
\toprule
Run period & Beam energy & Approx. EOT & Exposure used here & Prompt scan range \\
\midrule
2015 engineering run &
\SI{1.056}{GeV} &
$\sim 6.2\times10^{16}$ &
\SI{1170}{nb^{-1}} &
\SIrange{20}{130}{MeV} \\
2016 engineering run &
\SI{2.3}{GeV} &
$\sim 5.8\times10^{17}$ &
10\% development subset of a
\SI{10608}{nb^{-1}} parent sample &
\SIrange{35}{210}{MeV} \\
2021 physics run &
\SI{3.74}{GeV} &
$\sim 6\times10^{18}$ &
1\% review-stage subset of a $\sim\SI{160}{pb^{-1}}$ parent run &
\SIrange{30}{250}{MeV} \\
\bottomrule
\end{tabularx}
}
\caption{Data-taking campaigns carried through the present prompt-resonance note. The
quoted EOT values are campaign-scale approximations intended only to indicate the order
of magnitude of the accumulated statistics. In the present note, the fully unblinded
2015 engineering run anchors the methodology validation, the 2016 program is carried
through the current 10\% development sample on the way to the first 100\%
unblinding-stage analysis, and the 2021 program is represented by the current 1\%
validation subset on the way to the final 100\% analysis. The scan ranges
quoted here are the current GPR analysis ranges rather than the smaller published
window searches; the corresponding prompt-search publications scanned
\SIrange{19}{81}{MeV} in 2015 and \SIrange{39}{179}{MeV} in 2016
\cite{HPS2015DarkPhoton,HPS2016PromptLong}.}
\label{tab:dataset-summary}
\end{table}

The staged-analysis logic matters for the way the results are interpreted. The 2015
sample remains the cleanest environment for methodology validation because it is fully
unblinded. The 2016 and 2021 analyses are instead organized in review-friendly stages:
the present note uses the 2016 10\% and 2021 1\% subsets to validate the statistical
machinery, while the 2016 100\% first-unblinding milestone and the final fully
unblinded 2021 analysis remain the end points of the program. The 2021 1\% sample is
therefore best viewed as a validation subset of the upgraded detector configuration:
it already shows how the upgraded geometry and normalization inputs enter the
shared-$\eps^2$ reach, but it is not yet the final physics statement for the 2021 run.
